% questions.tex —— 1958 年·黑龙江、吉林、辽宁卷 · 题目内容片段（无答案）
% 由 tools/generate_projects.py 依来源总账自 1958/1958_黑龙江吉林辽宁_数学真题_试题.md 生成，勿手改（改源 MD 后重跑）


\begin{gktip}
\textbf{来源说明}：上游资料显示 1958 年黑龙江、吉林、辽宁三省高等学校招生考试同用这一份数学卷，上游源文件以三省并列命名，本卷按此收录。固定 PDF 共 1 页，页面为电子重排版（非原卷扫描件）；页面未印考试说明与逐题分值，题号 1—8 连续无缺跳，全卷无插图。
\end{gktip}

\begin{question}
\gkti{1}已知：圆 $O$ 内直交的二弦 $AB$、$CD$ 相交于 $P$，求证：$PA^{2}+PB^{2}+PC^{2}+PD^{2}$ 为一定值。

\end{question}
\begin{question}
\gkti{2}工地上有沙堆，其下底面为矩形，长宽各为 $a$ 米、$b$ 米，其各侧面都与地平面夹成 $45^\circ$ 角，上、下底距离为 $h$ 米，求沙堆体积。

\end{question}
\begin{question}
\gkti{3}求根式：$\sqrt[3]{a}$ 与 $\sqrt{5-a}$ 的乘积。

\end{question}
\begin{question}
\gkti{4}解方程：$x^{1+\lg x}=20^{\log_{20}100}$。

\end{question}
\begin{question}
\gkti{5}解方程组：

\nopagebreak[4]
\[
\begin{cases}
2x-5z=10-6a\\
x-2y+3z=2a+1\\
5x-6y+9z=23-10a
\end{cases}
\]

若 $x$ 为某直角三角形的斜边，$y$ 及 $z$ 为直角边，试确定 $a$ 的值。

\end{question}
\begin{question}
\gkti{6}已知：$\alpha$、$\beta$ 皆为锐角，求证：$\sin(\alpha+\beta)\le\sin\alpha+\sin\beta$。

\end{question}
\begin{question}
\gkti{7}已知：正方形 $ABCD$ 边长为 $a$，连接相邻各边中点，得一正方形 $A_{1}B_{1}C_{1}D_{1}$，再连接 $A_{1}B_{1}C_{1}D_{1}$ 相邻各边中点又得一个正方形 $A_{2}B_{2}C_{2}D_{2}$，如此相继下去，可得无数多个正方形。求这些正方形面积的总和。

\end{question}
\begin{question}
\gkti{8}解不等式：$2\cos^{2}x-(2+\sqrt3)\cos x+\sqrt3<0$。

\end{question}